% Source PDF page 17 - split at semantic boundaries
\begin{frame}[t]{Parameters (6/7)}
\BodyText{\textcolor{CommandGreen}{\texttt{-j}} tells iptables which target to use when a packet matches.}
\begin{itemize}
  \item Common targets: \textcolor{WarningColor}{\bfseries ACCEPT}, \textcolor{WarningColor}{\bfseries DROP}, \textcolor{WarningColor}{\bfseries LOG}, and \textcolor{WarningColor}{\bfseries REJECT}
  \item A rule may jump to a user-defined chain so additional rules can be applied
  \item Without a target, processing continues and the rule counter still increases
\end{itemize}
\end{frame}

\begin{frame}[t]{Parameters (7/7)}
\BodyText{\textcolor{CommandGreen}{\texttt{-o}} selects the outgoing interface.}
\begin{itemize}
  \item Filter table: valid with \textcolor{WarningColor}{\bfseries OUTPUT} and \textcolor{WarningColor}{\bfseries FORWARD}
  \item NAT/mangle tables: valid with \textcolor{WarningColor}{\bfseries POSTROUTING}
  \item Supports the same special options as \texttt{-i}
\end{itemize}
\end{frame}
